% 图契约：保留原稿的物理与电路拓扑，启用中文字体及中文断行，逐图检查标注。
\documentclass[tikz,border=14pt]{standalone}
\usepackage{fontspec,xeCJK}
\usepackage{amsmath,amssymb}
\usepackage{xcolor}
\usepackage{tikz}
\usetikzlibrary{arrows.meta,calc,patterns,positioning,decorations.pathmorphing,bending}
\IfFontExistsTF{Heiti SC}{\IfFontExistsTF{Heiti SC}{\setmainfont{Heiti SC}}{\setmainfont{FandolHei-Regular}}\IfFontExistsTF{Heiti SC}{\setsansfont{Heiti SC}}{\setsansfont{TeX Gyre Heros}}}{\IfFontExistsTF{Songti SC}{\setmainfont{Songti SC}}{\setmainfont{FandolSong-Regular}}\IfFontExistsTF{Songti SC}{\setsansfont{Songti SC}}{\setsansfont{FandolSong-Regular}}}
\definecolor{bluefield}{RGB}{24,76,255}
\definecolor{redfield}{RGB}{232,49,35}
\definecolor{orangefield}{RGB}{236,132,18}
\definecolor{greenfield}{RGB}{18,132,86}
\definecolor{purplefield}{RGB}{112,72,190}
\definecolor{linegray}{RGB}{60,68,78}
\definecolor{softblue}{RGB}{238,243,255}
\definecolor{softgreen}{RGB}{238,249,244}
\definecolor{softorange}{RGB}{255,247,235}
\definecolor{softpurple}{RGB}{245,240,255}
\tikzset{
  arrow/.style={-{Latex[length=3.1mm,width=2.1mm]},line width=1.1pt},
  thinarr/.style={-{Latex[length=2.5mm,width=1.7mm]},line width=.85pt},
  axis/.style={-{Latex[length=2.5mm,width=1.7mm]},line width=.8pt,draw=linegray},
  wire/.style={draw=linegray,line width=1.05pt},
  guide/.style={draw=linegray!55,dashed,line width=.65pt},
  box/.style={draw=linegray,rounded corners=3pt,fill=white,align=center,inner sep=6pt,line width=.75pt},
  bluebox/.style={box,draw=bluefield,fill=softblue},
  greenbox/.style={box,draw=greenfield,fill=softgreen},
  orangebox/.style={box,draw=orangefield,fill=softorange},
  purplebox/.style={box,draw=purplefield,fill=softpurple}
}
\IfFontExistsTF{Heiti SC}{\setCJKmainfont[AutoFakeBold=1.5]{Heiti SC}}{\setCJKmainfont[AutoFakeBold=1.5]{FandolHei-Regular}}
\IfFontExistsTF{Heiti SC}{\setCJKsansfont{Heiti SC}}{\setCJKsansfont{FandolHei-Regular}}
\begin{document}

\begin{tikzpicture}[x=1cm,y=1cm]
\path[use as bounding box] (-10,-4.6) rectangle (10,4.6);
\fill[white] (-10,-4.6) rectangle (10,4.6);
\node[font=\Large\bfseries,text width=18.6cm,align=center] at (0,4.05) {不同介质的电流：先辨认载流子和产生机制};
\begin{scope}[shift={(-5.4,-.25)}]
  \foreach \y/\lab/\col in {2.0/金属/bluefield,.7/电解质/greenfield,-.6/气体/orangefield,-1.9/真空/purplefield}{
    \node[font=\bfseries,\col] at (-3.0,\y) {\lab};
    \draw[linegray] (-1.9,\y)--(2.2,\y);
  }
  \foreach \x in {-1.4,-.7,0,.7,1.4}{\fill[bluefield] (\x,2) circle(.08);}
  \foreach \x in {-1.2,-.2,.8}{\node[greenfield] at (\x,.7) {$+$}; \node[orangefield] at ({\x+.45},.7) {$-$};}
  \fill[orangefield] (-.7,-.6) circle(.08); \fill[orangefield] (.5,-.6) circle(.08);
  \draw[arrow,bluefield] (-1.5,-.6)--(1.7,-.6);
  \draw[arrow,purplefield] (-1.5,-1.9)--(1.7,-1.9);
  \node[text=linegray] at (0,-2.65) {先辨认载流子，再写 $\mathbf J=\sum n_sq_s\mathbf v_s$};
  \end{scope}

\begin{scope}[shift={(4.55,-.15)}]
  \node[bluebox,text width=7.0cm,minimum height=2.15cm] at (0,1.25) {$\mathbf J=\sum_s n_sq_s\mathbf v_s$};
  \node[greenbox,text width=7.0cm,minimum height=2.15cm] at (0,-1.55) {负电荷逆场漂移，\\但对常规电流仍作同向贡献};
\end{scope}
\end{tikzpicture}
\end{document}
